```latex
\chapter{Yêu cầu giao diện bên ngoài}

Chương này mô tả các giao diện bên ngoài của hệ thống, bao gồm giao diện người dùng (User Interface), giao diện phần cứng, giao diện phần mềm và giao diện truyền thông. Các yêu cầu này xác định cách hệ thống tương tác với người dùng, các thành phần phần cứng và các hệ thống bên ngoài.

%%%%%%%%%%%%%%%%%%%%%%%%%%%%%%%%%%%%%%%%%%%%%%%%%%
\section{Giao diện người dùng (User Interface)}

Hệ thống được xây dựng dưới dạng ứng dụng Web với giao diện quản trị sử dụng FlexAdmin Dashboard Template kết hợp Bootstrap 5. Giao diện được thiết kế theo tiêu chí hiện đại, trực quan, dễ sử dụng và đảm bảo tính nhất quán giữa các chức năng.

Các thành phần chính của giao diện bao gồm:

\begin{itemize}
    \item Thanh điều hướng bên trái (Sidebar) dùng để truy cập các module của hệ thống.
    \item Thanh công cụ phía trên (Top Navigation Bar) hiển thị thông tin người dùng và các chức năng nhanh.
    \item Khu vực nội dung chính (Content Area) hiển thị dữ liệu của từng chức năng.
    \item Footer hiển thị thông tin phiên bản và bản quyền hệ thống.
\end{itemize}

Các yêu cầu đối với giao diện người dùng:

\begin{itemize}
    \item Giao diện phải thống nhất trên toàn bộ hệ thống.
    \item Hỗ trợ Responsive Design trên máy tính bảng và điện thoại.
    \item Các bảng dữ liệu hỗ trợ tìm kiếm (Search), sắp xếp (Sort), lọc dữ liệu (Filter) và phân trang (Pagination).
    \item Các biểu mẫu (Form) hỗ trợ kiểm tra dữ liệu đầu vào và hiển thị thông báo lỗi rõ ràng.
    \item Dashboard sử dụng biểu đồ trực quan để hiển thị dữ liệu thống kê.
    \item Màu sắc chủ đạo của giao diện sử dụng tông xanh lam và trắng nhằm phù hợp với nhận diện thương hiệu của doanh nghiệp.
\end{itemize}

Hình \ref{fig:ui_overview} minh họa bố cục tổng thể giao diện của hệ thống.

\begin{figure}[H]
    \centering
    \includegraphics[width=0.95\textwidth]{figures/ui/ui_overview.png}
    \caption{Bố cục tổng thể giao diện hệ thống}
    \label{fig:ui_overview}
\end{figure}

Hình \ref{fig:sitemap} mô tả cấu trúc điều hướng của hệ thống.

\begin{figure}[H]
    \centering
    \includegraphics[width=0.90\textwidth]{figures/ui/sitemap.png}
    \caption{Sơ đồ điều hướng các chức năng}
    \label{fig:sitemap}
\end{figure}

%%%%%%%%%%%%%%%%%%%%%%%%%%%%%%%%%%%%%%%%%%%%%%%%%%
\section{Giao diện phần cứng (Hardware Interface)}

Phiên bản hiện tại của hệ thống là Web Demo và chưa giao tiếp trực tiếp với thiết bị phần cứng. Tuy nhiên, kiến trúc được thiết kế để có thể tích hợp với hệ thống IoT trong tương lai.

Các thiết bị dự kiến kết nối với hệ thống bao gồm:

\begin{itemize}
    \item ESP32 Microcontroller.
    \item Cảm biến TDS.
    \item Cảm biến lưu lượng nước.
    \item Cảm biến nhiệt độ nước.
    \item Module Wi-Fi tích hợp trên ESP32.
\end{itemize}

Dữ liệu từ các cảm biến sẽ được gửi đến Server thông qua MQTT Broker để phục vụ việc theo dõi trạng thái hoạt động của thiết bị.

Hình \ref{fig:hardware_interface} minh họa mô hình kết nối phần cứng.

\begin{figure}[H]
    \centering
    \includegraphics[width=0.9\textwidth]{figures/architecture/hardware_architecture.png}
    \caption{Kiến trúc kết nối phần cứng}
    \label{fig:hardware_interface}
\end{figure}

%%%%%%%%%%%%%%%%%%%%%%%%%%%%%%%%%%%%%%%%%%%%%%%%%%
\section{Giao diện phần mềm (Software Interface)}

Hệ thống bao gồm các thành phần phần mềm chính như giao diện người dùng, ứng dụng PHP và hệ quản trị cơ sở dữ liệu MySQL.

Bảng \ref{tab:software_interface} mô tả các thành phần phần mềm được sử dụng.

\begin{table}[H]
\centering
\caption{Các thành phần phần mềm}
\label{tab:software_interface}

\begin{tabularx}{\textwidth}{|l|l|X|}
\hline
\textbf{Thành phần} &
\textbf{Công nghệ} &
\textbf{Chức năng} \\
\hline

Frontend &
HTML5, CSS3, JavaScript &
Xây dựng giao diện người dùng. \\
\hline

UI Framework &
Bootstrap 5 &
Thiết kế giao diện Responsive. \\
\hline

Admin Template &
FlexAdmin &
Cung cấp giao diện Dashboard quản trị. \\
\hline

Backend &
PHP &
Xử lý nghiệp vụ và điều phối dữ liệu. \\
\hline

Database &
MySQL &
Lưu trữ và quản lý dữ liệu hệ thống. \\
\hline

Web Server &
Apache (XAMPP) &
Triển khai và vận hành ứng dụng Web. \\
\hline

Version Control &
Git, GitHub &
Quản lý mã nguồn. \\
\hline

\end{tabularx}

\end{table}

Hình \ref{fig:software_architecture} mô tả kiến trúc phần mềm của hệ thống.

\begin{figure}[H]
    \centering
    \includegraphics[width=0.95\textwidth]{figures/architecture/software_architecture.png}
    \caption{Kiến trúc phần mềm}
    \label{fig:software_architecture}
\end{figure}

%%%%%%%%%%%%%%%%%%%%%%%%%%%%%%%%%%%%%%%%%%%%%%%%%%
\section{Giao diện truyền thông (Communication Interface)}

Hệ thống sử dụng các giao thức truyền thông phổ biến để trao đổi dữ liệu giữa các thành phần.

Các giao thức bao gồm:

\begin{itemize}
    \item HTTP/HTTPS giữa trình duyệt và Web Server.
    \item SQL giữa ứng dụng PHP và MySQL.
    \item MQTT giữa thiết bị IoT và MQTT Broker (dự kiến trong tương lai).
\end{itemize}

Quy trình truyền dữ liệu của hệ thống được mô tả như sau:

\begin{enumerate}
    \item Người dùng truy cập hệ thống thông qua trình duyệt Web.
    \item Trình duyệt gửi yêu cầu HTTP đến Web Server.
    \item Ứng dụng PHP xử lý yêu cầu nghiệp vụ.
    \item PHP thực hiện truy vấn hoặc cập nhật dữ liệu trong MySQL.
    \item Kết quả được trả về giao diện người dùng.
    \item Đối với dữ liệu thiết bị, ESP32 gửi dữ liệu đến MQTT Broker và hệ thống sẽ lấy dữ liệu để hiển thị trên Device Dashboard.
\end{enumerate}

Hình \ref{fig:communication_interface} minh họa luồng giao tiếp của hệ thống.

\begin{figure}[H]
    \centering
    \includegraphics[width=\textwidth]{figures/architecture/communication_architecture.png}
    \caption{Luồng giao tiếp giữa các thành phần}
    \label{fig:communication_interface}
\end{figure}
```
